\documentclass[12pt, a4paper]{article}
\usepackage{amsmath, amssymb, amsthm}
\usepackage{booktabs}
\usepackage{geometry}
\usepackage{hyperref}
\usepackage{enumitem}
\usepackage{xcolor}
\usepackage{mdframed}

\geometry{margin=2.5cm}

% Marker for illegible/missing handwritten content
\newcommand{\illegible}[1][?]{\textcolor{red}{[\textbf{ILLEGIBLE: #1}]}}

% Theorem environments
\newtheorem{theorem}{Theorem}[section]
\newtheorem{corollary}[theorem]{Corollary}
\theoremstyle{definition}
\newtheorem{definition}[theorem]{Definition}
\newtheorem{example}[theorem]{Example}
\theoremstyle{remark}
\newtheorem*{remark}{Remark}
\newtheorem*{note}{Note}

\title{\textbf{Series and Integral Transforms}\\[0.5em]
\large Class Notes Summary --- Weeks 1--3\\[0.3em]
\normalsize Sources: lecture notes + solved exercises from class sessions}
\author{Course 21183 --- HIT, Prof.\ Michael Kroyter}
\date{March 2026}

\begin{document}
\maketitle

\begin{mdframed}
\textbf{Key:} \illegible[description] = handwritten content that could not be read;
filled by mathematical context where possible.
\end{mdframed}

\tableofcontents
\newpage

% ==========================================================================
\section{Positive Series --- Convergence Tests, Cont.\ (18.3)}
% Source: מבחני התכנסות לטורים חיוביים - המשך, 18.3
% ==========================================================================

\subsection{Integral Test and Error Estimation (typed content)}

From the integral test, the sum of the generalized harmonic series is bounded:
\[
  \int_1^\infty f(x)\,dx \;\leq\; \sum_{n=1}^\infty a_n \;\leq\; a_1 + \int_1^\infty f(x)\,dx.
\]

\textbf{Case $p=2$, $f(x)=1/x^2$:} We obtain $1 \leq \sum 1/n^2 \leq 2$.
Using the partial sum $S_3 = 1 + \tfrac{1}{4} + \tfrac{1}{9} \approx 1.3611$ and the tail integral:
\[
  S_3 + \frac{1}{4} \;\leq\; \sum_{n=1}^\infty \frac{1}{n^2} \;\leq\; S_3 + a_4 + \frac{1}{4},
  \qquad \text{giving } 1.6111 \leq \sum \frac{1}{n^2} \leq 1.6736.
\]
(True value: $\pi^2/6 \approx 1.6449$.)

\begin{example}[$\sum \frac{1}{n \ln n}$ diverges --- typed]
$f(x) = 1/(x\ln x)$, substitution $u = \ln x$, $du = dx/x$:
\[
  \int \frac{dx}{x\ln x} = \int \frac{du}{u} = \ln(\ln x) \to \infty.
\]
So the series diverges.
\end{example}

\begin{example}[$\sum_{n=2}^\infty \frac{1}{n \ln n \ln(\ln n) \cdots}$]
Repeated substitution gives divergence for any finite chain of iterated logarithms.
(General principle: \illegible[precise statement of iterated log divergence criteria].)
\end{example}

\subsection{Ratio and Root Tests (18.3 --- mixed typed/handwritten)}

\begin{definition}
\[
  L = \limsup_{n\to\infty} \sqrt[n]{a_n} \quad \text{(root test / Cauchy)}, \qquad
  L = \lim_{n\to\infty} \frac{a_{n+1}}{a_n} \quad \text{(ratio test / D'Alembert)}.
\]
\end{definition}

\begin{theorem}
If $L < 1$: converges. If $L > 1$: diverges. ($0 \leq L \leq \infty$.)
\end{theorem}

\textbf{Key relation (typed):} If the ratio limit exists, then the root limit also exists and equals the same value. The converse is false --- root test is strictly stronger.

\textbf{Proof sketch for convergence ($L < 1$) --- typed:}
Set $q = (L+1)/2 < 1$. There exists $N$ such that for all $n > N$, $a_{n+1}/a_n < q$.
Then $a_{N+k} < q^k \cdot a_N$, so $\sum a_n$ is bounded by a geometric series. \qed

\textbf{Proof sketch for divergence ($L > 1$) --- typed:}
There exists $N$ such that $a_{n+1}/a_n > 1$ for all $n > N$, so $a_n \not\to 0$. \qed

\begin{example}[Interleaved sequence --- handwritten, partially legible]
\[
  \{a_n\} = \left\{1, 1, \tfrac{1}{2}, \tfrac{1}{3}, \tfrac{1}{4}, \tfrac{1}{9}, \tfrac{1}{8}, \tfrac{1}{27}, \ldots\right\},
  \quad
  a_n = \begin{cases} 1/2^{(n-1)/2} & n \text{ odd} \\ 1/3^{(n-2)/2} & n \text{ even.} \end{cases}
\]
Ratio test for odd $n$: $a_{n+1}/a_n = \frac{1/3^{(n-1)/2}}{1/2^{(n-1)/2}} = (2/3)^{(n-1)/2} \cdot \illegible[factor] \to 0$.\\
Ratio test for even $n$: $a_{n+1}/a_n = \frac{1/2^{n/2}}{1/3^{(n-2)/2}} = (3/2)^{(n-2)/2} \cdot \illegible[factor] \to \infty$.\\
So the ratio oscillates; generalized ratio test fails (even as a positive series!).

Root test:
\[
  \sqrt[n]{a_n} = \begin{cases} 2^{-(n-1)/(2n)} \to 1/\sqrt{2} & n \text{ odd} \\ 3^{-(n-2)/(2n)} \to 1/\sqrt{3} & n \text{ even.} \end{cases}
\]
Both limits are $< 1$, so $\limsup \sqrt[n]{a_n} = 1/\sqrt{2} < 1$. Converges.
\end{example}

% ==========================================================================
\section{Positive Series --- Completion (22.3)}
% Source: מבחני התכנסות לטורים חיוביים - סיום, 22.3
% ==========================================================================

\begin{example}[$\sum x^n/n^2$ --- handwritten]
Ratio test: $a_{n+1}/a_n = x \cdot n^2/(n+1)^2 \to x$.
\begin{itemize}
  \item $|x| < 1$: converges.
  \item $|x| > 1$: diverges.
  \item $x = 1$: $\sum 1/n^2$ converges (p-series).
  \item $x = -1$: \illegible[check at x=-1; likely converges conditionally].
\end{itemize}
Root test gives the same: $\sqrt[n]{x^n/n^2} = x \cdot n^{-2/n} \to x$.
\end{example}

\subsection{Cauchy Sequences (22.3 --- typed)}

\begin{definition}[Cauchy Sequence]
$\{a_n\}$ is Cauchy if $\;\forall \varepsilon > 0 \;\exists N : \forall n,m > N,\; |a_n - a_m| < \varepsilon$.
\end{definition}

\begin{theorem}[Cauchy's Criterion]
A sequence converges $\iff$ it is Cauchy.
\end{theorem}

Proof sketch (typed):
\begin{itemize}
  \item $(\Rightarrow)$: $|a_n - a_m| \leq |a_n - L| + |L - a_m| < 2\varepsilon$.
  \item $(\Leftarrow)$: Cauchy $\Rightarrow$ bounded $\Rightarrow$ convergent subsequence (Bolzano--Weierstrass) $\Rightarrow$ full sequence converges.
\end{itemize}

Why define this? (typed)
\begin{enumerate}
  \item Convergence detected without knowing the limit.
  \item Generalizes to metric spaces.
  \item Defines $\mathbb{R}$ as equivalence classes of Cauchy sequences of rationals.
\end{enumerate}

% ==========================================================================
\section{Absolute Convergence (25.3)}
% Source: התכנסות בהחלט, 25.3 (mixed typed/handwritten)
% ==========================================================================

\textbf{Setup (typed):} The only situation where positive-series tests cannot directly apply is when a series has infinitely many positive \emph{and} infinitely many negative terms.

\begin{definition}
\begin{itemize}
  \item $\sum a_n$ is \textbf{absolutely convergent} if $\sum |a_n| < \infty$.
  \item $\sum a_n$ is \textbf{conditionally convergent} if it converges but $\sum |a_n| = \infty$.
\end{itemize}
\end{definition}

\begin{theorem}
Absolutely convergent $\Rightarrow$ convergent.
\end{theorem}

Proof sketch (typed): $\sum |a_n|$ converges $\Rightarrow$ its partial sums are Cauchy $\Rightarrow$ partial sums of $\sum a_n$ are Cauchy (triangle inequality).

\begin{example}[Typed]
$\displaystyle\sum_{n=0}^\infty \frac{(-1)^n}{(2n)!}$: the absolute-value series is $\sum \frac{1}{(2n)!} = \cosh(1) \approx 1.543$.
By ratio test: $|a_{n+1}/a_n| = 1/((2n+1)(2n+2)) \to 0$. Converges absolutely.
Sum $= \cos(1) \approx 0.540$.
\end{example}

\begin{example}[Ratio/Root for general series --- partially handwritten]
Consider $a_n = \begin{cases} 5/7^n & n \text{ odd} \\ -1/7^n & n \text{ even} \end{cases}$, i.e.\ $\sum \frac{2+3(-1)^n}{7^n}$.

Ratio: $|a_{n+1}/a_n|$ alternates between $|-1/5|$ and $|-5|$ --- oscillates, inconclusive.

Root: $\sqrt[n]{|a_n|}$. For odd $n$: $(5/7^n)^{1/n} = 5^{1/n}/7 \to 1/7$.
For even $n$: $(1/7^n)^{1/n} = 1/7$.
So $\limsup \sqrt[n]{|a_n|} = 1/7 < 1$. Converges absolutely.
\end{example}

\begin{example}[$\sum x^n/n^2$ (absolute convergence domain) --- handwritten]
$|a_{n+1}/a_n| = |x| \cdot n^2/(n+1)^2 \to |x|$.
Absolutely convergent for $|x| < 1$.
At $|x| = 1$: $\sum 1/n^2$ converges, so absolutely convergent on the whole boundary too.
For $|x| > 1$: diverges (terms $\not\to 0$).
\end{example}

% ==========================================================================
\section{Alternating Series (29.3)}
% Source: טורים מחליפי סימן, 29.3 (mixed typed/handwritten)
% ==========================================================================

\subsection{Leibniz's Theorem (typed)}

A series $\sum a_n$ is called \textbf{alternating} if $a_n = (-1)^n b_n$ with $b_n \geq 0$.

\begin{theorem}[Leibniz]
If $b_n \geq 0$ and $b_n \searrow 0$ monotonically, then $\sum_{n=1}^\infty (-1)^{n+1} b_n$ converges.
Its sum $S$ satisfies $S_{2n} \leq S \leq S_{2n+1}$ and $|S - S_n| \leq b_{n+1}$.
\end{theorem}

\begin{proof}[Proof (handwritten, reconstructed)]
Even partial sums:
\[
  S_{2n} = (b_1-b_2) + (b_3-b_4) + \cdots + (b_{2n-1}-b_{2n}) \geq 0.
\]
Each bracket is $\geq 0$ (since $b_n$ decreasing), so $\{S_{2n}\}$ is increasing.
Also $S_{2n} = b_1 - (b_2-b_3) - \cdots - (b_{2n-2}-b_{2n-1}) - b_{2n} \leq b_1$, so bounded.

Odd partial sums:
\[
  S_{2n+1} = b_1 - (b_2-b_3) - (b_4-b_5) - \cdots - (b_{2n}-b_{2n+1}) \leq b_1.
\]
$\{S_{2n+1}\}$ is decreasing. Since $S_{2n+1} - S_{2n} = b_{2n+1} \to 0$, both subsequences converge to the same limit $S$.
Therefore $\{S_n\}$ converges.
\end{proof}

\begin{remark}
A Leibniz series cannot distinguish absolute from conditional convergence on its own.
\end{remark}

\begin{example}[Alternating harmonic series (typed)]
$\sum_{n=1}^\infty (-1)^{n+1}/n = 1 - 1/2 + 1/3 - 1/4 + \cdots = \ln 2 \approx 0.693$.

Partial sums: $S_1=1$, $S_2=1/2$, $S_3=5/6$, $S_4=7/12\approx0.583$, $S_5\approx0.783$, $S_6\approx0.617$.

Numerical acceleration (typed):
\[
  \frac{S_4+S_5}{2} \approx 0.683, \qquad \frac{S_5+S_6}{2} \approx 0.700.
\]
Exact: $S = \ln 2 \approx 0.693$.
\end{example}

\begin{example}[$\sum (-1)^n/n!$ (handwritten)]
$\sum 1/n! = e$, so $\sum (-1)^n/n!$ converges absolutely. Sum $= e^{-1} \approx 0.368$.
Partial sums $S_0=1, S_1=0, S_2=1/2, S_3=1/3, S_4\approx0.375$, $S_5\approx0.367$, $S_6\approx0.368$.
Bounds: $0.367 \leq S \leq 0.368$. \checkmark
\end{example}

\begin{example}[$\sum (-1)^{n+1}(n-2)^2/n^3$ --- handwritten]
$b_n = (n-2)^2/n^3$. Check $b_n \to 0$: yes. Check decreasing:
$b_n - b_{n+1} = \frac{n}{2n-1} - \frac{n+1}{2n+1} = \frac{1}{4n^2-1} > 0$. Yes.
Converges by Leibniz. Not absolute: $b_n \sim 1/n$, compare $\sum 1/n$ diverges.
\end{example}

\begin{example}[$\sum (-1)^{n+1} \frac{n}{2n-1}$ --- typed]
$b_n = n/(2n-1) \to 1/2 \neq 0$. Necessary condition fails. Diverges.
\end{example}

\subsection{Algebraic Properties of Series (typed)}

\begin{enumerate}
  \item $\sum c a_n = c \sum a_n$ and $c\sum a_n + c\sum b_n = c\sum(a_n+b_n)$.
  \item \textbf{Grouping}: inserting parentheses (same order) gives a grouping. \\
        If series converges $\Rightarrow$ every grouping converges to same sum. \\
        Converse fails (e.g.\ $\sum(-1)^n$). \\
        Exception: if all terms in each group have the same sign, grouping converging $\Rightarrow$ original converges.
  \item Changing finitely many terms does not affect convergence.
\end{enumerate}

\begin{example}[Grouping does not imply convergence --- typed]
$\sum_{n=0}^\infty(-1)^n$: grouping gives $(1-1)+(1-1)+\cdots = 0$, but partial sums $\{1,0,1,0,\ldots\}$ diverge.
\end{example}

\subsection{Rearrangements (typed)}

\begin{theorem}[Cauchy]
Rearranging terms of an \emph{absolutely} convergent series preserves the sum.
\end{theorem}

\begin{theorem}[Riemann Rearrangement]
If $\sum a_n$ converges conditionally, for any $L \in \mathbb{R}\cup\{\pm\infty\}$ there is a rearrangement with sum $L$.
\end{theorem}

\begin{proof}[Sketch (typed)]
Define $p_n = \max(a_n,0)$ and $N_n = \min(a_n,0)$. Since the series is conditionally convergent:
\[
  \sum |a_n| = \sum(p_n + N_n) = \infty, \quad \sum p_n = +\infty, \quad \sum N_n = -\infty.
\]
\textbf{To reach $L$}: greedily take positives until exceeding $L$, then negatives until below $L$; repeat. Overshoot $\to 0$ since $a_n \to 0$.\\
\textbf{To diverge to $+\infty$}: take enough positives to exceed $1$, one negative, enough positives to exceed $2$, etc.\\
\textbf{``Wild'' divergence}: alternate exceeding $+1$ then $-1$ \illegible[precise block construction].
\end{proof}

\begin{example}[Rearranging to $2\ln 2$ --- handwritten, partially legible]
\[
  \sum_{n=1}^\infty \frac{(-1)^{n+1}}{n} \cdot 2
  = 1 + \frac{1}{2} - \frac{1}{2} + \frac{1}{3} + \frac{1}{4} - \frac{1}{4} + \cdots
\]
\illegible[exact grouping/block argument for $2\ln 2$; blocked partial sums approach shown but not fully legible]
\end{example}

% ==========================================================================
\section{Solved Exercises from Tutorial Sessions}
% Sources: תרגול - דף 1 המשך (18.3), תרגול 3 - דף 2 (25.3)
% ==========================================================================

\subsection{SIT1 --- Tutorial Session Continuation (18.3)}
% From: תרגול - דף 1, המשך, 18.3

\textit{The following are exercises from Sheet 1 solved in the 18.3 tutorial session.}

\bigskip
\noindent\textbf{Exercise 1.5(ז):} $\displaystyle\sum_{n=1}^\infty \frac{1}{(2n+1)!}$

\textit{Solution:} Ratio test:
\[
  \frac{a_{n+1}}{a_n} = \frac{(2n+1)!}{(2n+3)!} = \frac{1}{(2n+2)(2n+3)} \to 0 < 1.
\]
Converges.

\bigskip
\noindent\textbf{Exercise 1.5(ו):} $\displaystyle\sum_{n=1}^\infty \frac{2n-1}{3^n}$

\textit{Solution:} Root test:
\[
  \sqrt[n]{\frac{2n-1}{3^n}} = \frac{\sqrt[n]{2n-1}}{3} \to \frac{1}{3} < 1.
\]
Converges.

\bigskip
\noindent\textbf{Exercise 1.5(א):} $\displaystyle\sum_{n=1}^\infty \frac{1}{2^n}\left(1+\frac{1}{n}\right)^{n^2}$

\textit{Solution:} Root test:
\[
  \sqrt[n]{\frac{1}{2^n}\left(1+\frac{1}{n}\right)^{n^2}} = \frac{1}{2}\cdot\left(1+\frac{1}{n}\right)^n \to \frac{e}{2} > 1.
\]
Diverges.

\bigskip
\noindent\textbf{Exercise 1.6(א):} $\displaystyle\sum_{n=0}^\infty \left(\frac{2}{3}\right)^n$

\textit{Solution:} Geometric series, $a_0 = 1$, $q = 2/3$. Sum $= \frac{1}{1-2/3} = 3$.

\bigskip
\noindent\textbf{Exercise 1.6(ז):} $\displaystyle\sum_{n=0}^\infty \left(\frac{7}{3^n} - \frac{6}{(n+3)(n+4)}\right)$

\textit{Solution:}
Geometric part: $\sum \frac{7}{3^n} = \frac{7}{1-1/3} = \frac{21}{2}$.

Telescoping part with $b_n = \frac{1}{n+3}$:
\[
  \sum_{n=0}^\infty \frac{6}{(n+3)(n+4)} = 6\sum_{n=0}^\infty\left(\frac{1}{n+3}-\frac{1}{n+4}\right) = 6\cdot\frac{1}{3} = 2.
\]
Total: $\frac{21}{2} - 2 = \frac{17}{2}$.

\bigskip
\noindent\textbf{Exercise 1.7:} Investigate convergence of $\displaystyle\sum_{n=1}^\infty \sin(nx)$ for all $x \in \mathbb{R}$.

\textit{Solution (handwritten, partially legible):}
For $x = k\pi$, $k\in\mathbb{Z}$: $\sin(nx)=0$ for all $n$. Converges trivially to 0.

For $x \neq k\pi$: Use the addition formula $\sin((n+1)x) = \sin(nx)\cos x + \sin x \cos(nx)$.
Since $\sin x \neq 0$, we have $\lim_{n\to\infty}\sin(nx)\cos(nx) = 0$ because \illegible[argument using $\sin^2+\cos^2=1$; claim that $\sin(nx)$ does not converge but partial sums are bounded].

The conclusion (typed at top of page): the series diverges for $x \neq k\pi$ \illegible[full justification illegible; likely via necessary condition since $\sin(nx)\not\to 0$, or via partial sums argument].

\bigskip
\noindent\textbf{Exercise 1.6(ה) [partially legible]:} \illegible[complex telescoping/geometric exercise involving $(-\sqrt{2})^n$ and $\sqrt[5]{n^2+n}$; full solution not recoverable]

\subsection{SIT2 --- Tutorial Session 3, Sheet 2 (25.3)}
% From: תרגול 3 - דף 2, 25.3

\textit{The following are exercises from Sheet 2 solved in the 25.3 tutorial session.}

\bigskip
\noindent\textbf{Exercise 2.2(א):} $\displaystyle\sum_{n=2}^\infty \frac{(-1)^n}{\ln n}$

\textit{Solution:}
$b_n = 1/\ln n \searrow 0$. By Leibniz: converges.
$\sum 1/\ln n > \sum 1/n$ (diverges by comparison), so not absolutely convergent.
\textbf{Conditionally convergent.}

\bigskip
\noindent\textbf{Exercise 2.2(ב):} $\displaystyle\sum_{n=1}^\infty \frac{(-1)^n}{\sqrt{n}}$

\textit{Solution:} $b_n = 1/\sqrt{n} \searrow 0$. Leibniz: converges.
$\sum 1/\sqrt{n}$ diverges (p-series, $p=1/2$). \textbf{Conditionally convergent.}

\bigskip
\noindent\textbf{Exercise 2.2(ד):} $\displaystyle\sum_{n=1}^\infty \frac{(-1)^{n+1}(n+1)}{n}$

\textit{Solution:} $b_n = (n+1)/n \to 1 \neq 0$. Necessary condition fails. \textbf{Diverges.}

\bigskip
\noindent\textbf{Exercise 2.2(ז):} $\displaystyle\sum_{n=1}^\infty \frac{(-1)^{n+1} n^3}{2^n}$

\textit{Solution:} Ratio test on $|a_n| = n^3/2^n$:
\[
  \left|\frac{a_{n+1}}{a_n}\right| = \frac{(n+1)^3}{2n^3} \cdot \frac{1}{2} \cdot \frac{2^n}{1}
  = \frac{1}{2}\left(\frac{n+1}{n}\right)^3 \to \frac{1}{2} < 1.
\]
\textbf{Absolutely convergent.}

\bigskip
\noindent\textbf{Exercise 2.3(ב):} $\displaystyle\sum_{n=1}^\infty \frac{1}{n\sin(1/n)}$

\textit{Solution:} $a_n = \frac{1}{n\sin(1/n)}$.
Since $\sin(1/n) \sim 1/n$ as $n\to\infty$, we get $a_n \to 1 \neq 0$.
Necessary condition fails. \textbf{Diverges.}

\bigskip
\noindent\textbf{Exercise 2.3(ד):} $\displaystyle\sum_{n=1}^\infty \frac{n+1}{n^2\sqrt{n+1}}$

\textit{Solution:} $a_n \sim 1/n^{3/2}$. Limit comparison with $1/n^{3/2}$:
$a_n \cdot n^{3/2} = n^{3/2}(n+1)/(n^2\sqrt{n+1}) \to 1$.
Since $\sum 1/n^{3/2}$ converges ($p=3/2>1$), the series \textbf{converges.}

\bigskip
\noindent\textbf{Exercise 2.3(ו):} $\illegible[series involving $n/e^n$; handwritten, illegible]$

\bigskip
\noindent\textbf{Exercise 2.4(א):} $\displaystyle\sum_{n=1}^\infty (-1)^n \frac{(2n)!}{4^n(n!)^2}$

\textit{Solution (handwritten):} Ratio test on $|a_n|$:
\[
  \left|\frac{a_{n+1}}{a_n}\right| = \frac{(2n+2)!}{4^{n+1}((n+1)!)^2}\cdot\frac{4^n(n!)^2}{(2n)!}
  = \frac{(2n+2)(2n+1)}{4(n+1)^2} = \frac{2n+1}{2(n+1)} \to 1.
\]
Ratio test inconclusive. \illegible[further analysis needed; likely shown not absolutely convergent by Stirling/Wallis, then Leibniz applied]

% ==========================================================================
\section{Office Hours --- Worked Proofs (23.3)}
% Source: שעת קבלה, 23.3
% ==========================================================================

\textit{The following material was covered in the 23.3 office hour session.}

\subsection{Integral Test --- Graphical Explanation (typed)}

The key diagram: for $f$ positive and decreasing with $f(n)=a_n$, the rectangle at $x\in[n,n+1]$ has area $a_n$. Comparing with the area under the curve:
\[
  \int_1^\infty f(x)\,dx \leq \sum_{n=1}^\infty a_n \leq a_1 + \int_1^\infty f(x)\,dx.
\]
The series and integral converge/diverge together.

\subsection{Limit Comparison Test --- Proof Sketch (handwritten)}

Given $b_n/a_n \to L$ with $0 < L < \infty$: for any $\varepsilon > 0$, eventually
\[
  (L-\varepsilon)\,a_n < b_n < (L+\varepsilon)\,a_n.
\]
So $\sum a_n$ and $\sum b_n$ are bounded by multiples of each other. They behave identically.

\subsection{Telescoping --- General Formula (handwritten)}

If $a_n = b_n - b_{n+1}$, then $S_n = b_1 - b_{n+1} \to b_1 - \lim b_n$.

Used in: $a_n = b_n - b_{n+1}$ where $b_n < a_n$; $\sum b_n$ and $\sum a_n$ controlled together when $b_n < a_n$ and $\sum b_n$ known.

\subsection{Exercise 1.8(a) --- Cauchy Criterion (handwritten)}

\noindent\textbf{Claim:} $\displaystyle\sum_{n=1}^\infty \frac{\cos(nx)-\cos((n+1)x)}{n}$ converges for all $x$.

\textit{Solution:} For $n > m$, the tail sum is:
\[
  \left|\sum_{k=m+1}^{n} a_k\right|
  = \left|\frac{\cos((m+1)x)}{m+1} - \frac{\cos((n+1)x)}{n} + \sum_{k=m+1}^{n}\cos((k+1)x)\left(\frac{1}{k+1}-\frac{1}{k}\right) + \cdots\right|
\]
After telescoping and bounding $|\cos| \leq 1$:
\[
  \leq \frac{1}{m+1} + \frac{1}{n} + \frac{1}{m+1} - \frac{1}{n+1} \leq \frac{3}{N} < \varepsilon,
\]
where $N = \lceil 3/\varepsilon \rceil$. Cauchy criterion holds for all $x$.

\subsection{Exercise 1.9(b) --- Divergence by Cauchy (handwritten)}

\noindent\textbf{Claim:} $\displaystyle\sum_{n=1}^\infty \frac{1}{\sqrt{n(n+1)}}$ diverges.

\textit{Solution:} Take $\varepsilon = 1/5$, $m = N$, $n = 2N$:
\[
  S_{2N} - S_N = \sum_{k=N+1}^{2N} \frac{1}{\sqrt{k(k+1)}}
  \geq \frac{N}{\sqrt{2N\cdot(2N+1)}} > \frac{N}{2N+1} > \frac{1}{5}.
\]
Cauchy criterion fails. Diverges.

\illegible[Additional office hour content --- several pages of handwritten derivations not fully legible; appeared to involve further convergence test examples]

% ==========================================================================
\section{Quick Reference: Convergence Test Decision Tree}
% ==========================================================================

\begin{enumerate}
  \item \textbf{Necessary condition}: $a_n \not\to 0$ $\Rightarrow$ diverges immediately.
  \item \textbf{Positive series}: try comparison, limit comparison, ratio, root, integral.
  \item \textbf{Alternating series}: check Leibniz ($b_n \searrow 0$) for conditional convergence; apply ratio/root to $|a_n|$ for absolute convergence.
  \item \textbf{General series}: apply ratio/root to $|a_n|$. If $L=1$, use Dirichlet or Abel.
  \item \textbf{When ratio test fails} (oscillating ratio): use generalized root test ($\limsup$).
  \item \textbf{Telescoping}: recognize $a_n = b_n - b_{n+1}$; sum $= b_1 - \lim b_n$.
  \item \textbf{Geometric}: $\sum a q^n$ converges iff $|q|<1$; sum $= a/(1-q)$.
\end{enumerate}

\end{document}
